{\scriptsize
    \setlength{\tabcolsep}{4pt}
    \renewcommand{\arraystretch}{1.10}
    \setlength{\LTleft}{0pt}
    \setlength{\LTright}{0pt}
    \setlength{\LTcapwidth}{\linewidth}
    \newlength{\PenelitianTerdahuluWidth}
    \setlength{\PenelitianTerdahuluWidth}{\dimexpr\linewidth-6\tabcolsep-3pt\relax}
    \begin{longtable}{@{}>{\raggedright\arraybackslash}p{0.14\PenelitianTerdahuluWidth}>{\raggedright\arraybackslash}p{0.30\PenelitianTerdahuluWidth}>{\raggedright\arraybackslash}p{0.18\PenelitianTerdahuluWidth}>{\raggedright\arraybackslash}p{0.38\PenelitianTerdahuluWidth}@{}}
        \caption{Penelitian Terdahulu yang Relevan}\label{tab:penelitian-terdahulu}                                                                                                                                                                                                                                                                                                                                                                                                                                                                                                                                                                             \\
        \toprule
        \textbf{Nama (Tahun)}                     & \textbf{Judul Penelitian}                                                                                                                                                            & \textbf{Metode Analisis}                                                                                                                       & \textbf{Hasil Penelitian}                                                                                                                                                                                                                                           \\
        \midrule
        \endfirsthead
        \toprule
        \textbf{Nama (Tahun)}                     & \textbf{Judul Penelitian}                                                                                                                                                            & \textbf{Metode Analisis}                                                                                                                       & \textbf{Hasil Penelitian}                                                                                                                                                                                                                                           \\
        \midrule
        \endhead
        \midrule
        \multicolumn{4}{r}{\footnotesize Dilanjutkan pada halaman berikutnya}                                                                                                                                                                                                                                                                                                                                                                                                                                                                                                                                                                                   \\
        \endfoot
        \bottomrule
        \endlastfoot
        \textcite{azhari2026gofood}               & Digital Marketing Strategy Based on E-Service, Brand Image and Reviews for GoFood MSMEs                                                                                              & Deskriptif kualitatif                                                                                                                          & Strategi digital pada UMKM GoFood perlu menekankan kualitas \textit{e-service}, penguatan \textit{brand image}, dan pengelolaan ulasan konsumen untuk meningkatkan daya saing dan penjualan.                                                                        \\
        \textcite{chen2026fashionanalytics}       & From Data to Decisions: The Role of Digital Marketing Analytics in Fashion E-commerce Performance                                                                                    & Kuantitatif; survei, PLS-SEM dengan SmartPLS 4.0, uji mediasi dan moderasi                                                                     & \textit{Digital marketing analytics} berpengaruh positif terhadap \textit{online retail performance} dan memediasi pengaruh sumber daya, kemampuan teknis, kemampuan manajerial, dan budaya berbasis data terhadap kinerja.                                         \\
        \textcite{portuguezcastro2026women}       & Enabling Women Entrepreneurs Through Digital Transformation: A Business Process Technology Perspective from Emerging Economies                                                       & \textit{Sequential exploratory mixed-method}; studi kualitatif dan survei pada 131 wirausaha perempuan                                         & Penelitian mengidentifikasi profil adopsi digital yang berbeda dan menekankan pentingnya pelatihan, dukungan teknologi proses bisnis, dan intervensi yang disesuaikan untuk memperkuat usaha perempuan di negara berkembang.                                        \\
        \textcite{touijer2026eo}                  & A Review of Entrepreneurial Orientation and Digital Transformation in SMEs: Toward a Multilevel Framework                                                                            & \textit{Systematic literature review} dengan prosedur PRISMA dan analisis bibliometrik atas 88 artikel                                         & Kajian ini menegaskan bahwa \textit{entrepreneurial orientation} menjadi pendorong penting transformasi digital UKM dan menghasilkan kerangka multilevel yang menghubungkan level individu, organisasi, dan lingkungan.                                             \\
        \textcite{zeren2026microenterprise}       & Social Media Marketing Adoption and Performance of Micro-Enterprise in an Emerging Market                                                                                            & Kuantitatif dengan Structural Equation Modeling (SEM) berbasis TOE                                                                             & Faktor teknologi, organisasi, dan lingkungan berpengaruh positif terhadap adopsi pemasaran media sosial, yang selanjutnya meningkatkan kinerja pemasaran dan performa usaha mikro.                                                                                  \\
        \textcite{agustin2025bumikayu}            & Strategi Pemasaran Digital Restoran Melalui Media Sosial dalam Keputusan Berkunjung (Studi pada Restoran Bumi Kayu di Kecamatan Samarang Kabupaten Garut)                            & Deskriptif kuantitatif                                                                                                                         & Pemasaran digital melalui media sosial, khususnya Instagram, berpengaruh positif terhadap keputusan berkunjung dan memperkuat pentingnya komunikasi visual restoran di platform digital.                                                                            \\
        \textcite{elverda2025ofdindonesia}        & Analysis of Factors Affecting Consumer Decision-Making in Choosing Online Food Delivery in Indonesia                                                                                 & Kuantitatif kausal dengan kuesioner, \textit{quota sampling}, dan SEM pada 724 responden                                                       & Persepsi harga, kegunaan, kemudahan penggunaan, kualitas layanan, kepercayaan, dan kepuasan memengaruhi niat beli ulang OFD; faktor paling dominan adalah \textit{perceived usefulness}.                                                                            \\
        \textcite{indriyani2025gudangbaju}        & Penerapan Strategi Pemasaran Digital Melalui Media Sosial dan E-Commerce dalam Meningkatkan Penjualan pada Gudang Baju Anak Lucu                                                     & Kualitatif deskriptif dengan studi kasus, wawancara, observasi digital, dan dokumentasi                                                        & Strategi melalui Instagram, TikTok, dan \textit{e-commerce} efektif meningkatkan visibilitas, \textit{engagement}, penjualan, \textit{repeat order}, dan rating toko, meski masih ada hambatan persaingan dan keterbatasan sumber daya.                             \\
        \textcite{ismail2025halalculinary}        & Strengthening the Global Competitiveness of Halal Culinary SMEs in North Sumatra through a Maqasid of Shariah-Based Approach                                                         & Kualitatif interpretif, \textit{multiple-case study}, wawancara mendalam, Atlas.ti 9, dan AHP                                                  & Permasalahan utama UMKM kuliner halal adalah \textit{digital marketing}, pengembangan produk, dan akses modal; strategi prioritas meliputi penguatan \textit{digital marketing}, peningkatan kualitas produk, pembiayaan syariah, dan percepatan sertifikasi halal. \\
        \textcite{kholifah2025tembindigo}         & Website-Based Digital Marketing and Digital Branding as a Strategy to Increase Ecoprint Product Sales at Tembindigo MSME                                                             & Participatory Action Research (PAR) dan diseminasi IPTEKS                                                                                      & \textit{Website}-based \textit{digital marketing} dan \textit{digital branding} meningkatkan eksposur bisnis, trafik pengunjung, dan penjualan; faktor pentingnya adalah konsistensi visual, kejelasan narasi, dan kemudahan pemesanan.                             \\
        \textcite{maharani2025kuliner}            & Strategi Pemasaran Digital untuk Meningkatkan Penjualan UMKM Sektor Kuliner                                                                                                          & Deskriptif kualitatif                                                                                                                          & Strategi digital melalui media sosial, katalog daring, dan pemesanan online meningkatkan visibilitas usaha, interaksi dengan konsumen, dan peluang kenaikan penjualan UMKM kuliner.                                                                                 \\
        \textcite{nurwansyah2025miegacoan}        & Customer Experience, E-Service Quality, and Loyality in Online Food Delivery: The Case of Mie Gacoan                                                                                 & Analisis deskriptif dan PLS-SEM pada 173 responden dengan \textit{voluntary sampling}                                                          & \textit{Customer experience} dan \textit{e-service quality} berpengaruh signifikan terhadap kepuasan, dan kepuasan menjadi mediator utama dalam pembentukan loyalitas pelanggan OFD Mie Gacoan.                                                                     \\
        \textcite{pradana2025infarm}              & Strategy for Utilizing TikTok Social Media as a Digital Marketing to Increase Sales at the Infarm Surabaya Brand                                                                     & Kualitatif deskriptif dengan analisis SWOT                                                                                                     & TikTok efektif meningkatkan \textit{awareness}, interaksi, dan peluang penjualan Infarm Surabaya ketika didukung oleh konten yang konsisten dan strategi digital yang terarah.                                                                                      \\
        \textcite{putro2025digitaltransformation} & The Influence of Digital Transformation On The Performance of Culinary SMEs in Bogor District Using The Technology Organization Environment (TOE) Model                              & Kuantitatif dengan model TOE dan pengujian hubungan antarvariabel                                                                              & Transformasi digital berpengaruh positif terhadap kinerja UMKM kuliner, dengan faktor teknologi, organisasi, dan lingkungan menjadi pendorong penting peningkatan performa usaha.                                                                                   \\
        \textcite{saptaji2025karawang}            & Digital Marketing Strategies for Household Processed Food MSMEs via Social Media and E-Commerce in Karawang                                                                          & Analisis deskriptif, IFE, EFE, IE, SWOT, dan AHP                                                                                               & Mayoritas UMKM telah memanfaatkan media sosial dan sebagian \textit{e-commerce}; hasil analisis menunjukkan kondisi mendukung strategi pertumbuhan dengan prioritas penguatan \textit{branding} dan penjualan digital terpadu.                                      \\
        \textcite{sophia2025rendanggadih}         & Analisis Strategi Pemasaran Rendang Gadih Melalui Platform Shopee untuk Meningkatkan Penjualan                                                                                       & Deskriptif kualitatif                                                                                                                          & Peningkatan penjualan dicapai melalui penerapan STP, bauran pemasaran 4P, serta pemanfaatan fitur Shopee seperti CRM, Live, Ads, dan Campaign.                                                                                                                      \\
        \textcite{trisaputra2025kamo}             & Efektivitas Promosi Konten Instagram Kamo Pet Care dan Implikasinya terhadap Strategi Pemasaran                                                                                      & Analisis efektivitas konten promosi Instagram dan implikasinya terhadap strategi pemasaran                                                     & Efektivitas konten Instagram terbukti penting untuk memperkuat komunikasi pemasaran, meningkatkan ketertarikan audiens, dan menjadi dasar perbaikan strategi promosi Kamo Pet Care.                                                                                 \\
        \textcite{azzahra2024guzzini}             & Marketing Strategy Analysis Using SWOT and QSPM Matrix (Case Study on Guzzini MSMEs)                                                                                                 & Kualitatif deskriptif; IFE, EFE, IE Matrix, SWOT Matrix, dan QSPM                                                                              & Posisi perusahaan berada pada sel \textit{grow and build}; SWOT menghasilkan enam alternatif strategi, dan prioritas QSPM menekankan promosi yang lebih kreatif serta pengaktifan toko offline dan online secara simultan.                                          \\
        \textcite{dini2024resilient}              & Formulating Strategies to Strengthen Resilient MSMEs in the F\&B Sector of Central Java: An Integrated SWOT, IFE, EFE, IE, and QSPM Approach                                         & \textit{Mixed-methods}; wawancara, kuesioner, IFE, EFE, IE Matrix, SWOT, dan QSPM                                                              & UMKM F\&B berada pada posisi \textit{grow and build}; strategi prioritas meliputi sertifikasi BPOM, pengamanan bahan baku, adopsi teknologi produksi, serta perluasan distribusi.                                                                                   \\
        \textcite{husein2024tagoya}               & Strategi Pemasaran Digital pada Usaha Mikro Kecil Menengah Tagoya Menggunakan Media TikTok                                                                                           & Analisis strategi pemasaran digital berbasis media TikTok                                                                                      & Pemanfaatan TikTok memperkuat visibilitas merek, memperluas jangkauan audiens, dan mendukung peningkatan penjualan UMKM Tagoya melalui konten yang lebih interaktif.                                                                                                \\
        \textcite{londono2024engagement}          & Engagement and Loyalty in Mobile Applications for Restaurant Home Deliveries                                                                                                         & PLS-SEM pada 349 pengguna aplikasi pesan antar restoran                                                                                        & \textit{Engagement}, kenyamanan, dan rasa aman berpengaruh positif terhadap loyalitas, sedangkan interaktivitas memengaruhi loyalitas secara tidak langsung melalui \textit{engagement}.                                                                            \\
        \textcite{pratama2024sambelmoji}          & Strategi Pemasaran Digital untuk Meningkatkan Penjualan UMKM Sambel Moji                                                                                                             & Analisis deskriptif                                                                                                                            & Solusi peningkatan penjualan difokuskan pada optimalisasi Instagram, pelatihan pemasaran digital, pemanfaatan \textit{tools} pemasaran, dan penguatan konten melalui \textit{influencer} atau \textit{content creator}.                                             \\
        \textcite{putri2024cikarang}              & Pemanfaatan Market Place GrabFood dan GoFood dalam Peningkatan Penjualan UMKM di Kecamatan Cikarang Barat                                                                            & Deskriptif kualitatif                                                                                                                          & Pemanfaatan GrabFood dan GoFood membantu UMKM memperluas jangkauan pasar, mempermudah transaksi, dan mendorong peningkatan penjualan di Kecamatan Cikarang Barat.                                                                                                   \\
        \textcite{sudirjo2024scarlett}            & Development of Promotional Criteria Decision Support System to Increase Sales of Scarlett Whitening Products through Digital Marketing Strategy                                      & Analisis kualitatif melalui wawancara dan simulasi untuk pengembangan sistem pendukung keputusan                                               & Strategi promosi digital, kolaborasi \textit{influencer}, dan sistem pendukung keputusan kriteria promosi dapat membantu meningkatkan efektivitas promosi dan penjualan produk Scarlett Whitening.                                                                  \\
        \textcite{zaheer2024ofda}                 & Enticing Attributes of Consumers' Purchase Intention to Use Online Food Delivery Applications (OFDAs) in a Developing Country                                                        & SEM-PLS, SmartPLS, \textit{bootstrapping}, dan CFA                                                                                             & \textit{Perceived value} berpengaruh positif terhadap \textit{purchase intention}, memediasi beberapa faktor utama, dan reputasi menjadi faktor dengan korelasi tertinggi terhadap niat beli pada aplikasi OFD.                                                     \\
        \textcite{haliza2023seqapla}              & Analisis Strategi Pemasaran dalam Upaya Peningkatan Penjualan pada Seqapla Coffee                                                                                                    & Analisis deskriptif kondisi internal-eksternal dan perumusan strategi pemasaran                                                                & Seqapla Coffee membutuhkan penguatan strategi pemasaran yang lebih terarah untuk menghadapi persaingan \textit{coffee shop} dan mendorong pemulihan penjualan.                                                                                                      \\
        \textcite{kong2023premiumcigars}          & The Promotion of Premium Cigars on Social Media                                                                                                                                      & Analisis isi kualitatif dan kuantitatif pada unggahan media sosial 47 merek cerutu                                                             & Media sosial dimanfaatkan untuk membangun citra gaya hidup, komunitas, dan loyalitas pelanggan melalui konten visual, informasi acara, serta interaksi langsung.                                                                                                    \\
        \textcite{kurniawan2023lamme}             & Pengembangan Strategi Pemasaran Potensi Gampong Lamme Melalui Media Digital                                                                                                          & \textit{Participatory research}; asesmen partisipatif, perencanaan, pelaksanaan, refleksi, evaluasi, implementasi marketplace dan media sosial & Pengembangan marketplace dan media sosial memperluas jangkauan pasar, meningkatkan potensi penjualan, serta mendukung penguatan ekonomi masyarakat melalui kolaborasi BUMG dan pelaku usaha.                                                                        \\
        \textcite{purnomo2023conversion}          & Digital Marketing Strategy to Increase Sales Conversion on E-commerce Platforms                                                                                                      & Deskriptif kuantitatif                                                                                                                         & Personalization, pemanfaatan media sosial, dan analisis kompetitor terbukti membantu meningkatkan konversi penjualan pada platform \textit{e-commerce}.                                                                                                             \\
        \textcite{rachmawati2023candisari}        & Implementation of Digital Marketing Strategy in MSME Development in Candisari Ungaran Village                                                                                        & Pendampingan implementasi digital marketing pada UMKM melalui observasi, sosialisasi, dan praktik                                              & Implementasi digital marketing membantu UMKM meningkatkan pengetahuan pemasaran, memperluas jangkauan pasar, dan membuka peluang peningkatan penjualan melalui pemanfaatan media sosial serta kanal digital.                                                        \\
        \textcite{ruhiyat2023minatbeli}           & Strategi Pemasaran Online untuk Meningkatkan Minat Beli Konsumen Produk Makanan dan Minuman UMKM di Kota Bogor                                                                       & Regresi linear berganda, analisis deskriptif, SWOT, IFE-EFE, IE, dan QSPM                                                                      & Unsur 4P berpengaruh positif terhadap minat beli; posisi UMKM berada pada strategi \textit{grow and build} dengan rekomendasi inovasi produk sesuai kebutuhan dan daya beli konsumen.                                                                               \\
        \textcite{santoso2023esatisfaction}       & The Role of E-Satisfaction on Repurchase and E-Wom Intention on The Costumers of Food Products By Local Micro and Small Businesses on The Digital Platforms                          & Kuantitatif dengan SEM-PLS                                                                                                                     & \textit{E-satisfaction} berpengaruh positif terhadap niat beli ulang dan \textit{e-WOM}, sehingga kepuasan digital menjadi faktor penting untuk mempertahankan pelanggan UMKM makanan di platform daring.                                                           \\
        \textcite{fatmawatie2022turnover}         & Digital Marketing Strategy to Increase Sales Turnover During the COVID-19 Pandemic                                                                                                   & Kajian deskriptif kualitatif                                                                                                                   & Digital marketing menjadi strategi adaptif berbiaya relatif rendah yang membantu pelaku usaha menjaga keberlangsungan bisnis dan meningkatkan omzet penjualan selama pandemi.                                                                                       \\
        \textcite{firdaus2022bangkupawon}         & Penerapan Strategi Digital Marketing pada Penggunaan Platform GoFood, GrabFood dan ShopeeFood terhadap Peningkatan Pendapatan Bisnis Kuliner (Studi Kasus Bangku Pawon Ngaliyan)     & Studi kasus deskriptif pada pemanfaatan platform OFD                                                                                           & Penggunaan GoFood, GrabFood, dan ShopeeFood terbukti membantu memperluas pasar dan meningkatkan pendapatan bisnis kuliner Bangku Pawon Ngaliyan.                                                                                                                    \\
        \textcite{karin2022instagram}             & Strategy for Using Instagram as a Digital Marketing Communication Media to Increase MSME Product Sales                                                                               & Deskriptif kualitatif                                                                                                                          & Instagram efektif digunakan sebagai media komunikasi pemasaran digital untuk memperluas jangkauan pasar dan meningkatkan penjualan produk UMKM.                                                                                                                     \\
        \textcite{khairani2022buahsayur}          & Strategi Meningkatkan Penjualan Buah dan Sayur pada Platform Digital di Kota Bogor                                                                                                   & Analisis deskriptif, regresi data panel, dan SWOT                                                                                              & Harga, ulasan, peringkat, jenis toko, dan promosi memengaruhi penjualan; strategi yang direkomendasikan adalah kerja sama langsung dengan produsen sebagai pemasok utama.                                                                                           \\
        \textcite{matos2022sodabrand}             & Global Case Study of Digital Marketing on Social Media by a Top Soda Brand                                                                                                           & Studi kasus global dengan analisis deskriptif pada aktivitas Facebook di 149 negara                                                            & Strategi pemasaran digital global bergantung pada frekuensi unggahan dan tingkat \textit{engagement}; konten visual dan interaksi yang konsisten meningkatkan popularitas merek, terutama di negara berkembang.                                                     \\
        \textcite{naufal2022raceofd}              & Perancangan Strategi Digital Marketing dengan Metode RACE pada Layanan Online Food Delivery Berdasarkan Perilaku Pelanggan Generasi Z Studi Kasus pada Go-Food dan GrabFood Surabaya & Metode RACE berbasis perilaku pelanggan Generasi Z                                                                                             & Strategi digital marketing yang disusun dengan kerangka RACE membantu memetakan tindakan \textit{reach}, \textit{act}, \textit{convert}, dan \textit{engage} sesuai perilaku pengguna OFD Generasi Z.                                                               \\
        \textcite{putri2022grabfood}              & Pengaruh E-service Quality dan E-trust terhadap E-satisfaction pada Pengguna Grab-Food Dari Generasi Z                                                                               & Kuantitatif eksplanatori dengan kuesioner dan analisis regresi berganda                                                                        & \textit{E-service quality} dan \textit{e-trust} berpengaruh positif signifikan terhadap \textit{e-satisfaction} pengguna GrabFood dari Generasi Z.                                                                                                                  \\
        \textcite{putri2022survivepandemic}       & Digital Marketing Strategy to Survive During COVID-19 Pandemic                                                                                                                       & Studi kasus kualitatif                                                                                                                         & BC Street Coffee bertahan selama pandemi melalui \textit{social media marketing} dan \textit{content marketing}, termasuk integrasi dengan GoFood dan GrabFood, sehingga penjualan daring menjadi lebih stabil dan meningkat.                                       \\
        \textcite{strydom2022churn}               & What Social Media Sentiment Tells Us About Why Customers Churn                                                                                                                       & Analisis sentimen media sosial selama 12 bulan terhadap lebih dari 1,7 juta unggahan                                                           & Terdapat tujuh pendorong utama churn, yaitu waktu tanggap, perilaku tidak etis, tagihan atau pembayaran, interaksi telepon, cabang atau toko, fraud atau scam, dan ketidakresponsifan.                                                                              \\
        \textcite{hidayah2021digitalplatforms}    & Effectiveness of Digital Platforms as Food and Beverage Marketing Media During the COVID-19 Pandemic                                                                                 & Wawancara dan studi literatur                                                                                                                  & Media sosial dan layanan \textit{online food delivery} efektif meningkatkan jangkauan pemasaran, volume penjualan, kesadaran merek, dan efisiensi operasional bisnis makanan dan minuman selama pandemi.                                                            \\
        \textcite{khotimah2021paklin}             & Analisis Strategi Digital Marketing dalam Meningkatkan Volume Penjualan Studi Kasus UMKM Paklin Donuts Kota Tegal                                                                    & Deskriptif kualitatif                                                                                                                          & Hasil penelitian menunjukkan bahwa digital marketing membantu meningkatkan volume penjualan UMKM Paklin Donuts Kota Tegal melalui promosi yang lebih luas dan komunikasi yang lebih cepat dengan konsumen.                                                          \\
        \textcite{mamuriyah2021gadihminang}       & Strategi Pemasaran dalam Meningkatkan Hasil Penjualan Rumah Makan Padang Gadih Minang                                                                                                & Kegiatan PKM dengan promosi digital melalui brosur, Instagram, dan YouTube                                                                     & Promosi melalui brosur, Instagram, dan YouTube diarahkan untuk memperkenalkan Rumah Makan Padang Gadih Minang ke khalayak yang lebih luas dan diharapkan mendorong kenaikan omzet penjualan bulanan.                                                                \\
        \textcite{purnama2021ecommerceumkm}       & Analisis E-commerce dalam Membantu Penjualan UMKM di Tengah Pandemi                                                                                                                  & Kualitatif deskriptif melalui studi literatur dan tinjauan sumber data sekunder                                                                & \textit{E-commerce} menjadi solusi penting bagi UMKM untuk memperluas pasar dan menjaga efisiensi operasional pada saat mobilitas konsumen menurun.                                                                                                                 \\
        \textcite{roslam2021markisaaurora}        & Strategi Pemasaran UMKM Markisa Aurora di Kota Makassar pada Masa Pandemi COVID-19                                                                                                   & Analisis deskriptif strategi pemasaran UMKM pada masa pandemi                                                                                  & UMKM Markisa Aurora dapat mempertahankan dan memperluas pasar melalui penguatan promosi digital, penyesuaian strategi pemasaran, dan pemanfaatan kanal penjualan daring selama pandemi.                                                                             \\
        \textcite{suharjo2020bogorsnacks}         & Strategi Pemasaran Digital pada UMKM Makanan Ringan di Kota Bogor                                                                                                                    & Analisis deskriptif, perilaku pembelian daring konsumen, \textit{Importance Performance Analysis}, bauran pemasaran digital, dan 4C            & Strategi prioritas menekankan penguatan konten digital, pemanfaatan media sosial, dan penyesuaian elemen bauran pemasaran digital untuk memperluas jangkauan pasar UMKM makanan ringan di Kota Bogor.                                                               \\
        \textcite{gurel2017swot}                  & SWOT Analysis: A Theoretical Review                                                                                                                                                  & Tinjauan teoretis terhadap konsep dan penggunaan SWOT                                                                                          & Kajian ini menegaskan bahwa SWOT merupakan alat analisis strategis yang membantu organisasi mengidentifikasi kekuatan, kelemahan, peluang, dan ancaman sebagai dasar perumusan strategi.                                                                            \\
        \textcite{helms2010swotreview}            & Exploring SWOT Analysis -- Where Are We Now?: A Review of Academic Research from the Last Decade                                                                                     & Kajian pustaka terhadap penelitian SWOT dalam satu dekade                                                                                      & Hasil telaah menunjukkan bahwa SWOT tetap luas digunakan dalam penelitian strategi, tetapi memerlukan penerapan yang lebih sistematis dan dukungan kerangka analitis yang lebih kuat agar hasilnya lebih operasional.                                               \\
    \end{longtable}
}
